\chapter{Discussion}
\label{chap:discussion}

The results indicate that chirps might play a role in mediating competitive
interactions among Brown Ghost Knifefish. By developing and implementing a deep
learning pipeline for automated chirp detection, I was able to observe and
analyze over 73,000 chirps during unconstrained interactions, surpassing any
previous study in similar conditions. These findings reveal that subordinate
fish emit chirps, which are correlated with the cessation of aggressive
behaviors. This suggests a communicative function beyond active sensing.

The introduction of this deep learning pipeline represents a significant
advancement in the methodology for studying electrocommunication, as this fully
automated approach enables the detailed analysis of large datasets and the
study of freely interacting animals in both laboratory and field settings. This
marks a step towards quantitative ethology of electrocommunication.

This following chapter will unfold in two segments: Initially, I will discuss
the implications and limitations of the behavioral correlates with the chirps
of \textit{A.~leptorhynchus}. The focus will then shift to the chirp detection
pipeline itself—evaluating its efficacy, possibilities for enhancement, and
applications in research.

\section{Are chirps a signal of submission?}

Chirps are produced in a variety of behavioral contexts. They can occur
spontaneously without any obvious external stimuli
\parencite{englerSpontaneousModulationsElectric2000a}. They can also be induced
by playing back a sine wave that mimics the electric organ discharge (EOD) of a
conspecific. Although the chirp rate decreases if the simulated conspecific
also produces chirps \parencite{walzNeuroethologyElectrocommunicationHow2013}.
Studies both in field and laboratory settings have also documented
echo-responses to chirping
\parencite{henningerStatisticsNaturalCommunication2018,
hupeEffectDifferenceFrequency2008, gamasalgadoEchoResponseChirping2011}. In the
context of mating, chirps have been observed both in the laboratory and the
field, suggesting a role in courtship and synchronization of spawning.
\textcite{hagedornCourtSparkElectric1985a} reported that a playback of male
chirps could trigger spawning in females, even without the presence of a real
male. \textcite{henningerStatisticsNaturalCommunication2018} demonstrated
stereotyped chirping patterns during courtship for two species, one recorded in
the laboratory and one in the field. They also provided observations of
chirps produced by retreating intruders that disturbed courting pairs.
Competitive interactions recorded in controlled conditions suggest opposing
functions of chirps. They have been correlated with a reduction in attack
rates, implying they might serve as a submissive signal
\parencite{hupeElectrocommunicationSpeciesWeakly2012}. In another cases,
dominant individuals were observed to produce more chirps, suggesting a role in
signaling dominance or aggression
\parencite{triefenbachChangesSignallingAgonistic2008b}. This inconsistency
along with the lack of clear behavioral correlates led
\textcite{obotiWhyBrownGhost2023} to propose that chirps might serve as probing
signals, a form of autocommunication to enhance sensory information during
social interactions.

In line with the findings of
\textcite{hupeElectrocommunicationSpeciesWeakly2012,
henningerStatisticsNaturalCommunication2018}, I have shown that in the context
of shelter competitions in \textit{A. leptorhynchus}, chirps are
disproportionally produced by fish that lost the competitions. For a signal to
be classified as communication, it needs to be a deliberate action, structure,
or trait produced by one animal (the sender) to convey information to another
animal (the receiver), influencing the receiver's behavior in a way that
typically benefits the sender
\parencite{bradburyPrinciplesAnimalCommunication2011, bergstromEvolution2016,
alcockAnimalBehavior2019}. As the chirps produced by the subordinate fish
strongly correlate with the offset of chasing events, similar to observations
by \textcite{henningerStatisticsNaturalCommunication2018}, we hypothesized that
chirps might be a means of communicating submission, and thus stopping the
chasing of the dominant. But in this case, chirps should not have correlated
precisely \emph{with} the offset of chasing, but with a point in time that is
shortly \emph{before} the offset. Additionally, ritualized chasing events
lasted longer when the loser fish produced chirps (although we cannot rule out
that the chasing events, in which chirps were produced, could have been longer
if no chirps were emitted). Because most of these chirps were produced
precisely at (or around) the end of chasing, it is sensible to assume that
prolonged chasing caused the subordinate to emit chirps, instead of the other
way around. Another explanation could be that chirps might not alter the
dominants behavior in the current competition, but increase the time it takes
for the dominant to attack again. But chirps produced by the submissive during
chasing showed no effect on the latency between attacks by the dominant.
Whether chirps of the submissive have an effect on the probability of getting
bit by the dominant has yet to be explored. Following the previous definition
of a communication signal, the reports of echo responses strongly suggest that
chirps serve as such because they change the behavior of the receiver
\parencite{henningerStatisticsNaturalCommunication2018,
hupeEffectDifferenceFrequency2008, gamasalgadoEchoResponseChirping2011}. But in
the context of the competitions analyzed here no such change in the behavior
of the receiver could be found. So while we observed clear correlations between
chirping and social status as well as with a specific phases of a ritualized
competition, it remains unclear what function chirps might serve in this
context.

The lack of evidence on the function of chirps as communication signals in this
setting does not automatically support the assumption that they do not serve a
communicative function at all. The correlations not only with social status but
also with specific spatio-temporal interactions stand in contrast with the
interpretation of \textcite{obotiWhyBrownGhost2023}, which argue that chirps
might rather be a means of active sensing through autocommunication that
enhance the perception of conspecifics. In the experiment of
\textcite{raabElectrocommunicationSignalsIndicate2021}, fish competed in
absolute darkness, leaving mainly the lateral line and electrolocation as
reliable sensory inputs. During chasing, fish stay close to each other for
prolonged times and under high velocities. The chasing fish and the one being
chased require temporally and spatially precise sensory inputs to coordinate
this movement. If chirps were emitted to perceive a conspecific, there is no
obvious explanation to why only the subordinate individual would require to do
so, specifically at the time where chasing ends. At the same time, the option
that chirps might explicitly transmit information to conspecifics does not rule
out \textcite{obotiWhyBrownGhost2023} proposal that chirps could be used for
active sensing. As suggested in other systems, such as in bats and whales,
signals that are used for active sensing could be used for communicating in
different contexts \parencite{mossAdaptiveEcholocationBehavior2023}.

The largest limitation of the competition arena by
\textcite{raabElectrocommunicationSignalsIndicate2021} is arguably that it is
confined in space. Observations of competitions in the field are constrained to
a few events observed by
\textcite{henningerStatisticsNaturalCommunication2018}, where an intruder male
interrupts a courting dyad. In response, the courting male interrupts
courtship, and approaches the intruder from a distance of more than a meter, to
which the intruder reacts with chirps and retreats. Retreat was not an option
in this study, as the fish were constrained to the aquarium. This could also be
an explanation to why competitions lasted the full three hours of the active
period in this setting. \textcite{raabElectrocommunicationSignalsIndicate2021}
showed that similar to chirps so-called frequency rises, another type of
electrocommunication signal evolving on a time scale of seconds, are
predominantly produced by subordinates and strongly correlate with the
\emph{onset} of competitions. This was interpreted in the way, that rises are
used by the subordinate to continue conspecific assessment by motivating the
dominant to compete. But while this temporal correlation of rises and chasing
onsets is strong, only 10\% of the chasing events were actually preceded by a
rise emitted by a subordinate. So the majority of competition events were
probably not initialized by electrocommunication signals but either by spatial
interactions or just spontaneously motivated by the sheer presence of another
fish in the perceivable range of the dominant individual. In the field,
territories of single fish could be larger than the experimental tank and a
scenario, in which fish are forced to compete for prolonged time, would not
occur. Following this assumption, chirps in combination with a retreat could
still signal submission in naturalistic conditions. The missing physical
possibility of retreat in the dataset could explain a lack of an effect of the
subordinate chirps onto the dominants behavior. Such multimodalities in
communication signals are extensively demonstrated for courtship signaling in
many systems and may also apply to competitions for shelter
\parencite{ronaldMaleMiceAdjust2020,
ronaldWhatMakesMultimodal2017,uetzComplexSignalsComparative2017,
redingContextdependentPreferencesVary2017,
mitoyenEvolutionFunctionMultimodal2019}. To further explore this possibility,
we either need to go back to the field to record competition scenarios in
naturalistic settings or stage competitions with the possibility of retreat.

\section{Automated detection of communication signals}

Deep learning tools for analyzing animal communication signals have been
demonstrated for the vocalizations of a wide variety of taxa including
primates, whales, rodents, bats, birds, frogs, and fish
\parencite{coffeyDeepSqueakDeepLearningbased2019,
cohenAutomatedAnnotationBirdsong2022, bellafkirBatEcholocationCall2022,
berglerORCASPOTAutomaticKiller2019, lawsonAutomatedAcousticDetection2023a,
kimuraEvaluationDeepLearningBased2023,
waddellApplyingArtificialIntelligence2021, erbsAutomatedLongtermAcoustic2023}.
In all previous cases, the models are trained entirely on tediously
hand-labeled data. The problems they solve can usually be classified into two
categories: The first one are species classifiers used on soundscape
recordings from the field, typically trained to classify spectrogram snippets
of animal vocalizations into different species during \acrfull{PAM}. These
systems often employ \acrshortpl{CNN} \parencite{kahlBirdNETDeepLearning2021,
kimuraEvaluationDeepLearningBased2023,
waddellApplyingArtificialIntelligence2021} or transformers
\parencite{bellafkirBatEcholocationCall2022}. The second category encompasses
song detectors, usually applied to some focal species in laboratory
experiments. Here, the focus is on locating the vocalizations in time (and
frequency) rather than identifying the species. Detection can be addressed
using standard CNNs \parencite{berglerORCASPOTAutomaticKiller2019}, but more
complex tasks, such as call segmentation and classification, often utilize RCNN
or YOLO models \parencite{coffeyDeepSqueakDeepLearningbased2019,
katherDevelopmentMachineLearning2024, lawsonAutomatedAcousticDetection2023a}.
An exception is TweetyNet, which incorporates \acrfull{LSTM} layers
\parencite{cohenAutomatedAnnotationBirdsong2022}.

All systems mentioned are limited to \emph{detection}. In electric fish, the
produced continuous electric field serves as an electric tag, enabling the
tracking of individuals over time \parencite{hagedornCourtSparkElectric1985a,
henningerTrackingActivityPatterns2020,
madhavHighresolutionBehavioralMapping2018,
raabAdvancesNoninvasiveTracking2022}. Since fish can be separated by frequency
and chirps are constrained to a specific relative frequency range, merely
detecting the time of chirps is insufficient. So I chose a model that predicts
chirp boundaries on the frequency axis, as previously demonstrated for rodents
\parencite{coffeyDeepSqueakDeepLearningbased2019}. Associating detected signals
with the emitting individual is a problem of \emph{assignment}. Because most
other system lack a robust passive tag, this problem remains unaddressed in any
of the presented applications. This thesis presents a first approach to emitter
assignment using feature engineering and a simple CNN for electric fish.
Together, the detection and assignment model form a continuous pipeline for
automatic chirp detection and assignment in large datasets. Instead of relying
on large hand-labeled training datasets, I additionally developed a simulation
routine to generate artificial recordings of electric fish. This way, I can
automatically generate training data from a small representative sample of
communication signals. This reduces the need for the error-prone process of
hand-labeling large amounts of training data.

\subsection{Performance of the detection and simulation pipeline}

In terms of performance, both the metrics used to evaluate the detection model
as well as the assignment step are promising: The average precision reaches
more that 85\% (AP@[.5:.01:.95], figure \ref{fig:detector_performance}) for
detection and 97\% for assignment (figure \ref{fig:assignment_training}). From
these metrics, detection and assignment errors are expected to be quite rare.
But it is important to emphasize that the detection performance is evaluated
partially on simulated data and the assignment model fully on simulated data,
which probably leads to overestimating the performance of the pipeline. For
this reason, I manually checked hundreds of spectrograms in addition to the
performance metrics, two of which are introduced in detail in figures
\ref{fig:chirp_detections_simultaneous} and
\ref{fig:chirp_detections_bigchirps}. These spectrograms illustrate that for
most cases and in recordings with an adequate \acrfullpl{SNR}, the detection
pipeline performs much better compared to previous implementations,
particularly for longer and larger chirps. First applications on field data
should paint a clearer picture on how usable this approach is and what needs to
be improved. A promising approach to further increase the performance of the
detection model could be through active learning. Here, the model labels
previously unseen data. From these detections, uncertain samples, such as those
with low confidence scores, are picked, manually annotated and the model is
retrained on those. This is a thoroughly tested method of increasing the
performance of machine learning algorithms
\parencite{zhanComparativeSurveyDeep2022}.

In addition to a thorough evaluation of the detection pipeline it is equally
important to evaluate and improve the simulation pipeline. Electric fish offer
the unique opportunity to capture their full behavior by electric recordings
but also to simulate a complete scene of multiple interacting fish. In such
scenarios, even a trained human observer struggles to assign fish to the
correct emitter. This makes the availability of \emph{true} ground-truths that
do not rely on human annotators indispensable. The current simulation routine
is merely a first proof of principle, where many aspects of natural electrode
grid recordings are still missing.

Currently, it starts by generating EOD$f$ traces that include communication
signals such as chirps and rises. Based on those frequency traces, the EOD is
generated from Fourier components to resemble that of \textit{A.
leptorhynchus}. A separate pathway generates movement of this simulated fish
and the EOD is then amplitude-modulated on each simulated electrode to reflect
the current distance between the fish and each electrode.

To improve this pipeline, we can start at the simulation of communication
signals. Chirps are currently generated by extracting the evolution of the
EOD$f$ from chirps of real fish and inserting them into the simulated EOD$f$ of
the simulated fish. This way, the diversity of the simulations of chirps is
constrained by the available data. Instead of using the extracted instantaneous
frequencies of chirps directly, we could train a learning
algorithm to generate artificial EOD$f$s of chirps. Algorithms such as
\acrfullpl{VAE} \parencite{nishikimiVariationalAutoencoderBased2024} or
\acrfullpl{GAN} \parencite{dissanayakeGeneralizedGenerativeDeep2023,
hazraSynSigGANGenerativeAdversarial2020} are already used in medical
applications to synthesize EEG, EMG and ECG signals and could be applied in a
similar fashion to synthesize the frequency evolution of the EOD during chirps
or rises. Instead of the real EOD$f$ traces, we could then use the synthetic
communication signals to simulate fish.

The second step would be the improvement of the simulation on the level of the
EOD itself. In the current simulation, the fish is an oscillating monopole and
the electric field decays equally in all directions. But in reality, the
electric field is an oscillating dipole with two lobes, which does not decay
equally in all directions \parencite{bendaPhysicsElectrosensoryWorlds2020}.
This difference would not be visible on the sum spectrograms I used for chirp
detection, which is why I ignored it for now. But when future versions of chirp
detectors incorporate spatial information into the assignment step, this
feature could become important. Multiple models of such an electric field
already exist and could be implemented and incorporated into the current
pipeline \parencite{babineauModelingElectricField2006,
nelsonPreyCaptureWeakly1999, wallachInternalModelCanceling2023}. As a result,
the way in which the field decays would not just depend on position, but on
body size and curvature of the fish.

After adding realistic simulation of the dipole, we should also generate the
fish movement based on the statistics of real movements. This is
important, because these motions substantially influence the amplitude
modulations of the EOD. These are visible on spectrograms and perceivable by
other fish \parencite{yuElectrosensoryContrastSignals2019,
fotowatStatisticsElectrosensoryInput2013a}. In the current simulation, only the
distribution of swimming velocity is taken from approximations from field
recordings. Other patterns, such as forward-backward motions and heading
directions, are hand-tuned only to \emph{resemble} those of a real fish.
Instead, we should use data from experiments where animals were video-tracked.
This would allow us to generate models that not only produce realistic
velocities but also trajectories, body curvature, and maybe even behavioral
patterns such as foraging, fighting or resting. All of those would bring us
closer to reliably recreate the full statistics of natural recordings and hence
closer to true ground truths to build and benchmark analysis algorithms such as
the chirp detector.

One of the most important aspects of the present simulation pipeline is the
augmentation with real noise. Adding realistic noise to simulated recordings is
not only essential to enable the transfer to real data but also provides
endless possibilities to \emph{augment} the simulations. The same simulation on
many different \enquote{noise backdrops} can drastically increase the number of
training samples for learning algorithms
\parencite{salamonScaperLibrarySoundscape2017}. Currently, this noise is taken
from another recording of fish in the laboratory. Instead, we should aim to
create a library of noise backdrops from different settings in the lab an in
the field. This would not require much effort, since the recording hardware
already exists and fish would not be needed. Not just the background noise
itself but also the \acrfull{SNR} is crucial. The current simulation pipeline
is restricted to a single SNR. Instead, we could use various levels of SNRs to
generate many training samples for a single simulation-noise combination, which
would make the model more robust to weak EOD amplitudes or fish that are far
away from the electrodes.

The last step would be to incorporate the complexity of the environment in such
simulations. Many objects substantially distort the electric field
\parencite{zlenkoVisualizationElectricFields2023,
pereiraContextualEffectsSmall2005, bacherNewMethodSimulation1983} and the
natural habitat of \textit{A. leptorhyncus} is a complex three dimensional
labyrinth of stacked rocks \parencite{hagedornCourtSparkElectric1985a}. The
surface structure of such habitats could be reconstructed using underwater
photogrammetry \parencite{franciscoHighresolutionNoninvasiveAnimal2020}. The
problem with this approach is that fish prefer to hide in crevices between
those rocks, which are usually not picked up well by conventional
photogrammetry approaches. How such a complex 3D environment could be included
in these simulations is yet to be explored.

Improving the simulation pipeline would also facilitate a comparison between
the chirp detector's and human-level performance, the gold-standard in
supervised detection and classification tasks. This way of evaluating
performance is hardly ever viable, because in most cases, the ground truth is
generated by human annotators. This way, a model logically never achieves to
surpass human level performance. But in a case where the actual ground truth is
accessible, such as in simulations, human level performance can be compared to
the detectors performance. Under noisy conditions, this could enable the
training of detection pipelines that could surpass the performance of human
annotators.

\subsection{Convolutional neural networks versus transformer architectures}

RCNN and YOLO architectures, such as the ones used in this pipeline, have
already been used successfully for the automatic detection of ultrasonic
vocalizations of rodents \parencite{coffeyDeepSqueakDeepLearningbased2019}.
Here, I provide another successful application of object detection models to
the detection of electric signals. But in the last years, the development of
CNNs for computer vision and audio analysis has slowed down due to the rise of
transformer-based architectures. Transformers were mainly developed for natural
language processing and became popular through the public demonstration of the
\acrfull{GPT}. Since then, transformer-based architectures have been
successfully applied to vision tasks and usually outperform convolution-only
architectures in terms of accuracy and precision
\parencite{shehzadiObjectDetectionTransformers2023}. But convolutional
approaches still surpass transformers, when the computational time is
considered. In the work presented here, the processing time had to be as low as
possible to facilitate the fast processing of weeks of recordings. But
transformers will likely become more performant in the future. The
computational costs will decrease and the available hardware will become
better. At that point, we should move to transformer models for
Neuroethological signal recognition.

\subsection{Back to the time domain}

Wave-type weakly electric fish offer the unique property that in opposite to
bats or whales, the signals they produce are continuously emitted. Furthermore,
each EOD approximately consists of just a single frequency at a point in time.
As mentioned in \hyperref[fig:multifish_chirps]{chapter two}, if we record a
single fish, we can estimate this frequency for a very fine temporal precision
by the inverse of each period. But as soon as noise is introduced or the
signals of multiple fish overlap, this becomes impossible as illustrated in
figure \ref{fig:multifish_chirps}. So
\textcite{raabAdvancesNoninvasiveTracking2022,
henningerTrackingActivityPatterns2020} developed the \texttt{wavetracker} to
approximate the frequency changes over time of multiple fish on spectrograms. I
followed with the \texttt{chirpdetector} that is capable of also detecting
chirps on spectrograms of multiple fish. But working on these spectrograms
comes with a large pitfall: We are always restricted by the time-frequency
uncertainty. As an example, the spectrograms used in the \texttt{chirpdetector}
have a temporal resolution of \SI{20}{\milli\second}, corresponding to a single
frequency estimate for short chirps. This is just enough to make them visible
but still reduces the frequency resolution in a way that fish less than
\SI{20}{\hertz} apart are hardly separable from each other. But what if we were
able to extract each fish from such a recording similarly to how we can record
single fish in the laboratory? I believe that this is not so far away as we
might think: Convolution and transformer architectures running on spectrograms,
waveforms or both simultaneously can be used for sound source separation,
meaning that they can separate e.g., vocals, bass and drums from a rock song
\parencite{rouardHybridTransformersMusic2022,
stollerWaveUNetMultiScaleNeural2018, petermannCocktailForkProblem2022}.
Applications of sound source separation in Biology have already been
demonstrated successfully for vocalizations of birds and fish
\parencite{dentonImprovingBirdClassification2022,
mancusiFishSoundsEvaluation2022}. Our problem is a bit different. In the case
of separating instruments from vocals (or bird songs of different species)  the
task is \emph{semantic} segmentation. If we would try to apply this approach to
electric fish, i.e., separating a recording of multiple fish into recordings of
single fish, we would be faced with \emph{instance} segmentation. Instead of
segmenting EODs from the background, the task is to segment the EODs of
individual fish from each other \emph{and} the background. While this sounds
complex, the fact that fish are usually separated in frequency could make this
approachable. But we would need a learning algorithm that combines methods from
sound source separation with methods from instance segmentation from computer
vision. An improved simulation pipeline would be an ideal way of training such
models. Having temporal resolutions not constrained to that of spectrograms
would be an enormous benefit that allows for much more intricate insights into
the social lives of electric fish. This is because we do not just use frequency
to understand their communications, but also amplitude to estimate their
positions.

\subsection{Applications of chirp detection}

With these technical considerations in mind, for what other practical
applications could the \texttt{chirpdetector} be useful? Electric fish are one
of the main fractions that make up the total biomass of fish in the Amazon
basin \parencite{lassoCompositionAbundanceBiomass1992,
cramptonEcologicalPerspectiveDiversity2011a}. They present a large variability
in habitat use, social lives and communication repertoires that are
comparatively poorly explored and still offer exiting discoveries to be made.
With the right training data, the chirp detector could be trained and used on
the chirps of any chirping species. With the simulation pipeline to
generate training data, just a small representative subset might be enough to
get good results. In combination with an electrode grid and the
\texttt{wavetracker}, we now have the tools to record and analyze electric and
spatio-temporal behaviors of multiple, freely interacting wave-type weakly
electric fish in their natural habitat.
\textcite{henningerStatisticsNaturalCommunication2018} demonstrated that this
approach is indispensable in reflecting upon decades of neuroethological
research in electric fish. \textcite{fortuneSpookyInteractionDistance2020a}
showed that this approach can be employed to study behavioral differences
between surface and cave dwelling populations of two closely
related species. With the pipeline presented here, we can add the quantitative
analysis of chirps to this suite of methods. Just by the sheer amount of chirps
that are produced in various behavioral contexts, without arguing for or
against a certain function, it is sensible to assume that they are
important to the lives of electric fish
\parencite{henningerStatisticsNaturalCommunication2018, obotiWhyBrownGhost2023,
hupeElectrocommunicationSignalsFree2009}. Because chirps are connected
tightly to courtship or spawning, they could be used as proxies for mating or
at least courtship events
\parencite{henningerStatisticsNaturalCommunication2018,
hagedornCourtSparkElectric1985a}. Currently not even the mating system of the
different species of wave-type electric fish is known. Beyond ethological and
comparative field studies, this approach is equally applicable to recordings of
multiple fish in controlled conditions. \textcite{obotiWhyBrownGhost2023}
correctly criticized the lack of causal evidence for the function of chirps as
a communication signal. The pipeline presented here should make it much easier
to study chirps and enable a much higher throughput compared to manual
annotations, to shed a light onto the function(s) of chirps.

In parallel to fundamental research, the chirps of weakly electric fish could
become equally important in the applied sciences, such as in population
monitoring. The electric sense is widely regarded as a significant evolutionary
boon, contributing substantially to the remarkable diversity and abundance of
electric fish \parencite{cramptonEvolutionElectricSignal2006}. This unique
adaptation, however, necessitates specific conditions in terms of nutritional
availability and water chemistry, rendering these species potentially
vulnerable to anthropogenic disturbances, as discussed by
\textcite{markhamEnergeticsSensingCommunication2016}. An aspect not covered in
their discussion is the potential exacerbation of these vulnerabilities by
electrical noise. In areas close to human settlements, recordings of electric
fish in their natural habitats are often marred by significant electrical
interference within their perceptible frequency ranges—a problem not
encountered in more remote locations. A particularly concerning source of such
interference is the proliferation of hydroelectric dams
\parencite{agostinhoAllThatGoes2011, keppelerEarlyImpactsLargest2022,
duponchelleConservationMigratoryFishes2021}. These structures not only pose
barriers to riverine migration but also generate substantial electrical noise.
Monitoring fish populations is essential to prevent further damage to their
diversity and abundance \parencite{larentisEffectsHumanDisturbance2022}. Active
monitoring by telemetry is possible but expensive, error prone and restricted
to few large species \parencite{hahnBiotelemetryRevealsMigratory2019}.
Monitoring through eDNA, where species accounts are extracted from the DNA
available in water samples, is a promising and holistic approach and removes
the need for more invasive monitoring methods
\parencite{desantanaCriticalRoleNatural2021}. To complement these approaches,
purely passive monitoring was demonstrated recently in the endangered Amazonian
River Dolphins, using hydrophones and neural networks to detect echolocation
signals \parencite{erbsAutomatedLongtermAcoustic2023}. But this is restricted
to two species. I would argue that we can take this a step further with
wave-type, and maybe at some point even pulse-type electric fish.
\textcite{henningerStatisticsNaturalCommunication2018} and
\textcite{madhavHighresolutionBehavioralMapping2018} have demonstrated that
electric fish can not only be passively monitored, but that we can even track
individual identities over time and estimate their positions to a comparatively
fine temporal scale. What has been missing were their communication signals. If
we would be able to monitor and automatically analyze the communication of
these fish, we would not only be able to asses how their abundances change over
time but also how the reproductive behavior of these populations is affected by
environmental changes. This could be implemented in a purely passive way,
without the need to catch fish or significantly disturb their environment.
 
\section{Conclusion} % (fold)
\label{sec:conclusion}

The automated chirp detection pipeline developed in this thesis successfully
detects all known types of chirps in recordings of multiple Brown Ghost
Knifefish. This represents a significant advancement in the study of
electrocommunication, since it is the first system capable of detecting and
automatically assigning each chirp to the fish that emitted it.

Applying this pipeline to a dataset from staged competitions between two fish
surpassed the number of chirps in most previous studies and revealed important
behavioral insights. Chirps were predominantly emitted by subordinate fish and
were correlated with the cessation of aggressive interactions by dominant
individuals. This finding suggests a communicative function of chirps in
mediating ritualized competitions for resources.

In conclusion, the chirp detection pipeline lays the groundwork for deeper
exploration into the electrocommunication of wave-type electric fish. It paves
the way for studying social dynamics, behavioral cues, and ecological roles of
chirping in both laboratory and natural settings and marks a first step towards
non-invasive and passive monitoring of electrocommunication in wild
populations.

% section Conclusion (end)

% chapter Discussion (end)

\iffalse
The disparity between the current capabilities in data analyses and the
practices with which data is not only analyzed but also recorded in some fields
of research is striking. While controlled experiments and
\enquote{conventional} statistics are indispensable, whether such causal
effects revealed in experiments play a role in the natural environment and
behavioral context of species is often overlooked in the push for causality in
strictly controlled conditions. But especially now, as not only the autonomous
recording of enormous amounts of data but also automated analyses become more
and more accessible, going into the field and inspiring new hypotheses by
quantitative observations becomes approachable. As an example, smartphone apps
can automatically recognize bird songs and classify the species just from
comparatively poor-quality phone recordings
\parencite{kahlBirdNETDeepLearning2021}. Despite this, many studies still rely
on manually labeling communication signals of all kinds of species. Software
such as
\href{http://www.mackenziemathislab.org/deeplabcut}{\texttt{DeepLabCut}} and
\href{https://sleap.ai/}{\texttt{SLEAP}} enables automated tracking of
individual body parts of most animals \parencite{pereiraSLEAPDeepLearning2022a,
lauerMultianimalPoseEstimation2022}. This provides the possibility of high
throughput in behavioral experiments and frees time that can be allocated to,
for example, recording more data. But computational technologies are still not
widely adopted nor included into teaching curricula. Instead, many rely on
cheap manual labor in the form of undergraduate students to annotate datasets
by hand. \textcite{andersonScienceComputationalEthology2014} nicely summarized
this process as beeing slow, imprecise, subjective, low-dimensional, limited by
the human sensory system and lastly \emph{dull}. The need for annotated
datasets is not something that we soon be able to completely resolve, but
instead of labeling a full dataset by hand, it is much less time consuming to
label a small representative subset to train a machine learning system that
then labels the full dataset for analysis. This would decouple the amount of
data that can be recorded from the number (or patience) of available
annotators. This way, junior scientists, often times still students themselves,
would not be prevented from learning skills that become more and more valuable
not only in science but already are so in industry.

\fi
